\section{Discussion}\label{sec:discussion}

\subsection{What heterogeneity costs, and what it does not}
The usual explanation for FedAvg's shortfall in household forecasting is that homes are too different. Our
controlled experiment does not support that explanation for final accuracy. Holding everything else fixed, doubling
the load-shape divergence between clients changed FedAvg's test error by at most 0.0012\,kWh, comparable to the
seed-to-seed variation.
The fragmentation sweep shows where the gap does come from. The federated protocol with one client matches
centralised training, and the error then grows with the logarithm of the number of clients the same data are spread
over, whether or not those clients differ.
Heterogeneity did matter for the speed of convergence: the traffic needed to approach centralised accuracy grew by
about 85\,\% from random to clustered grouping. In practice this means that heterogeneity should be budgeted as a
communication and latency cost. It is not a reason to expect a worse final model, provided that the round budget is
large enough and the best round is chosen on validation data.

This also explains why remedies aimed at client drift showed little benefit here. FedProx helped only when its
coefficient was close to zero, and the adaptive variant performed worse. When fragmentation, rather than drift between
heterogeneous clients, causes most of the gap, a drift penalty has little to correct and mainly slows local progress.

\subsection{Implications for deployment}
\begin{itemize}
\item \emph{Federate few, large clients.} The single most consequential design choice was the number of clients
the data were spread over. Aggregating homes behind a building gateway, a community energy hub or a supplier-side
aggregator and federating those (here $K=5$) cost \EonefourfedavgKfiveGap{}\,\% in MAE relative to centralised training, with a quarter
of the traffic of 20 clients (\EonefourfedavgKfiveMB{}\,MB), whereas meter-level federation ($K=100$) lost \EonethreefedavghouseholdGap{}\,\% and
moved about 30 times the raw data. Where governance permits an intermediate tier, it is the cheapest accuracy lever
available.
\item \emph{Federation is a governance choice, not a bandwidth saving, at this scale.} A household's eight months of training
half-hourly readings occupy about 140\,kB, only about twice the size of a single message of this small LSTM. Thirty rounds
of federated training therefore move several times more data than uploading the readings would, and even the
traffic to near-centralised accuracy is of the same order as the data. The case for FL in residential forecasting
rests on data minimisation, consent and regulatory constraints, not on network load. This agrees with the qualitative
argument of Briggs et al.~\cite{briggs2020ccai}, and it qualifies analytic estimates of large bandwidth
savings~\cite{taik2020,savi2021}, which depend on model size, sampling rate and the number of rounds.
\item \emph{Fine-tune locally after federating.} One epoch of local fine-tuning costs no communication. It closed
much of the gap between FedAvg and local training per household, and it reduced the share of homes harmed by
federation from around two thirds to below one half. Any federated household forecaster should ship with this step.
\item \emph{Select the round on validation error.} Validation error rose temporarily in early rounds in every
federated configuration. A fixed, short round budget can therefore deliver a worse model than a longer budget with
validation-based selection.
\item \emph{Watch the high-consumption tail.} The homes that gain most from pooled data are the highest consumers,
and they are also the ones FedAvg serves worst. A distribution system operator using federated forecasts for
congestion management should evaluate errors per consumption tier, not only on average.
\end{itemize}

\subsection{On adaptive proximal coefficients}
HA-FedProx was designed on the premise that the dispersion of client updates measures heterogeneity. The premise is
half right. The dispersion ratio does rank partitions by heterogeneity, but within a run it is dominated by the
convergence of the global model: as the shared descent direction shrinks, any residual disagreement, including pure
stochastic-gradient noise, makes the ratio approach one. A schedule proportional to $\rho_t$ therefore increases the
proximal pull late in training regardless of the data. Two corrections suggest themselves: normalising $\rho_t$ by
its value under a random partition, which requires a reference the server does not have, or using only early-round
dispersion, which our data show is the most heterogeneity-sensitive. We leave both to future work. Our main
recommendation is a caution: dissimilarity estimates computed from client updates (here, an update-space version of
FedProx's $B$-dissimilarity) confound heterogeneity with the stage of optimisation, and any rule that adapts
regularisation to them should be checked against this confound.

\subsection{Limitations}
\begin{itemize}
\item \emph{Scope.} One dataset, one city, one study year and one horizon: a winter test period and a 30-minute
horizon for 100 standard-tariff London homes. The conclusions concern a small LSTM. Larger models raise
per-message traffic proportionally, which strengthens the communication conclusion but may change the accuracy
results.
\item \emph{Simulated network and privacy.} The network and stragglers are simulated, and traffic counts payload
bytes only; protocol overhead was measured at about 1\,\%. Keeping data local is not a formal privacy guarantee.
Updates can leak information~\cite{zhu2019leakage,geiping2020inverting}, and neither differential privacy nor secure
aggregation was applied. Both would add accuracy or communication costs on top of those measured here.
\item \emph{Fixed hyper-parameters.} The learning rate, model size and window were fixed a priori and shared by all
arms. Only the proximal coefficients were tuned. A per-arm search could move individual numbers, although the
controlled comparison is designed to make such shifts affect all arms alike.
\item \emph{Server-side evaluation.} For efficiency the simulator scores the global model on every household's
validation and test windows. A deployment would compute these scores on the clients and report only the metrics.
\end{itemize}
